% Page 32 du document source
% Transcription provisoire — à relire et corriger si nécessaire.

\begin{center}
    {\Large\itshape Alban pendant la guerre de Cent Ans}
\end{center}

\vspace{1.5em}

En septembre 1392, Alban reçoit Hélias de Vern, procureur du chapitre
Sainte Cécile d'Albi. Les consuls albigeois envoient deux sergents le
rejoindre. Ils sont porteurs d'un acte de procédure. Ces personnages ne
font que transiter par Alban. Ils se rendent dans le bas Languedoc. Et
ceci prouve l'importance de cette voie
de grand trafic à cette époque, de la route Aix -- Montauban sur laquelle
Alban se trouve comme un point d'étape apprécié.

\medskip

En cette fin de siècle, la famille des « Nougaret » semble prendre quelque
importance chez nous. Un Jacques de Nougaret est « viguier » (juge) à Albi
de 1396 à 1405.

\medskip

Au début du XV\textsuperscript{ème} siècle (Jeanne d'Arc naît en 1412),
une autre famille d'Alban, les De Moret, semble prendre de l'importance
dans notre coin. À cette date, en effet, Jehan de Moret prête serment pour
un fief noble qu'il détient dans la juridiction d'Alban. Et vingt ans
après en 1433, un Dominique de Moret nous est signalé comme marchand et
notable de notre cité. En 1450, nous apprenons qu'il fait commerce de fer
-- de ce fer qu'on extrait des « ferrières » de Paulin. Ce qu'on appelle
encore la Mouline (où résidait jadis ce brave Louis Fages) comportait en
fait deux moulins : l'un « drapier » parce qu'on y devait faire de la
toile ; c'était la « Mouline Bergan », et l'autre au-delà du ruisseau
« la Mouline Ferran », où se faisait le traitement des minerais de fer…
Aujourd'hui, il ne reste plus que le souvenir de ce temps passé, puisque
la « Mouline » n'a plus d'habitants.
